% !TEX program = xelatex
% ---------------------------------------------------------------------------
% Simple Homework -- an academic worksheet.
% Everything above \begin{document} is styling. Edit the "Assignment details"
% block, then replace the example problems in the body. Lines tagged
% CUSTOMIZE are the usual places to change the look.
% ---------------------------------------------------------------------------
\documentclass[11pt,a4paper]{article}

% ---------- Packages ----------
\usepackage{fontspec}
\usepackage[a4paper,margin=2.5cm]{geometry}   % CUSTOMIZE: page margins
\usepackage{amsmath,amssymb,amsthm}
\usepackage{microtype}
\usepackage{xcolor}
\usepackage{parskip}
\usepackage{enumitem}
\usepackage{titlesec}
\usepackage{fancyhdr}
\usepackage{booktabs,tabularx,array}

% ---------- Fonts ----------
% TeX Gyre Termes is the primary font. TeX Gyre Heros is not in the bundled
% font catalogue, so labels and headings use Termes (bold, letter-spaced).
\setmainfont{TeX Gyre Termes}                 % CUSTOMIZE: body font

% ---------- Palette: neutrals plus one academic accent ----------
\definecolor{Ink}{HTML}{1C1C1C}      % body text
\definecolor{Muted}{HTML}{5C6066}    % running header, notes, small labels
\definecolor{Rule}{HTML}{CFCFCF}     % hairlines and answer-box borders
\definecolor{Shade}{HTML}{F3F3F1}    % instruction box background
\definecolor{Accent}{HTML}{1F3A5F}   % CUSTOMIZE: the single accent (navy)
\AtBeginDocument{\color{Ink}}

% ---------- Text and spacing ----------
\linespread{1.1}                                       % CUSTOMIZE: line spacing
\setlength{\parskip}{6pt plus 1.5pt minus 1pt}         % CUSTOMIZE: paragraph spacing
\setlength{\emergencystretch}{3em}
\setlength{\fboxsep}{8pt}
\renewcommand{\arraystretch}{1.25}
\setlength{\heavyrulewidth}{0.8pt}
\setlength{\lightrulewidth}{0.4pt}
\setcounter{secnumdepth}{0}                            % \section headings are unnumbered

% ---------- Running header and footer ----------
\setlength{\headheight}{14pt}
\pagestyle{fancy}
\fancyhf{}
\fancyhead[L]{\small\color{Muted}\coursename}
\fancyhead[R]{\small\color{Muted}\studentname}
\fancyfoot[C]{\small\color{Muted}\thepage}
\renewcommand{\headrulewidth}{0.4pt}
\renewcommand{\headrule}{{\color{Rule}\hrule height\headrulewidth width\headwidth}\vskip-\headrulewidth}
\fancypagestyle{plain}{%  first page: footer only
  \fancyhf{}\fancyfoot[C]{\small\color{Muted}\thepage}%
  \renewcommand{\headrulewidth}{0pt}}

% ---------- Section headings (use \section*{Part A -- ...}) ----------
\titleformat{\section}
  {\normalsize\bfseries\color{Accent}\addfontfeature{LetterSpace=8}}
  {}{0pt}{\MakeUppercase}
  [\vspace{3pt}{\color{Accent}\titlerule[0.8pt]}]
\titleformat{\subsection}{\large\bfseries\color{Ink}}{}{0pt}{}
\titlespacing*{\section}{0pt}{2.4ex plus 0.5ex minus 0.2ex}{1.4ex plus 0.2ex}
\titlespacing*{\subsection}{0pt}{2ex plus 0.4ex minus 0.2ex}{0.9ex plus 0.2ex}

% ---------- Lists: parts, nested parts and multiple choice ----------
\setlist{itemsep=0.3em, topsep=0.3em, parsep=0.3em, partopsep=0pt}
\setlist[itemize,1]{label=\textcolor{Accent}{\textbullet}, leftmargin=1.6em}
\setlist[enumerate,1]{label=\textcolor{Accent}{\bfseries\arabic*.}, leftmargin=1.8em}
% \begin{parts} \item ... \end{parts}  gives (a), (b), ...; nest it for (i), (ii), ...
\newlist{parts}{enumerate}{2}
\setlist[parts,1]{label=\textcolor{Accent}{\bfseries(\alph*)}, leftmargin=2.4em, itemsep=0.5em}
\setlist[parts,2]{label=\textcolor{Accent}{(\roman*)}, leftmargin=2.2em, itemsep=0.3em}
% \begin{choices} \item ... \end{choices}  gives A., B., C., D.
\newlist{choices}{enumerate}{1}
\setlist[choices,1]{label=\textcolor{Accent}{\bfseries\Alph*.}, leftmargin=2em,
  itemsep=0.15em, parsep=0pt, topsep=0.3em}

% ---------- Problems ----------
% \problem[points]{Topic}   (the points argument is optional)
\newcounter{problem}
\numberwithin{equation}{problem}   % equations read (3.1), (3.2), ... per problem
\newcommand{\problem}[2][]{%
  \refstepcounter{problem}%
  \par\addvspace{1.2\bigskipamount}%
  \noindent{\large\bfseries\color{Accent}Problem~\theproblem}%
  \enspace{\large\bfseries\color{Ink}#2}\hfill{\small\color{Muted}#1}\par\nobreak
  \vspace{-0.2em}\noindent{\color{Rule}\rule{\linewidth}{0.4pt}}\par\nobreak}

% ---------- Solution, explanation and note styles ----------
% \begin{solution}[Label] ... \end{solution}  (label defaults to "Solution")
\newenvironment{solution}[1][Solution]
  {\par\smallskip\noindent\textbf{\color{Accent}#1.}\enspace\ignorespaces}
  {\par\smallskip}
% \begin{note}[Label] ... \end{note}  (label defaults to "Note"; try [Hint])
\newenvironment{note}[1][Note]
  {\par\small\color{Muted}\noindent\textbf{#1.}\enspace\ignorespaces}
  {\par\smallskip}
\renewcommand{\qedsymbol}{\ensuremath{\blacksquare}}   % end-of-proof mark

% ---------- Instruction box (shaded) ----------
\newsavebox{\LPShadeBox}
\newenvironment{instructions}[1][Instructions]
  {\par\medskip\begin{lrbox}{\LPShadeBox}%
   \begin{minipage}{\dimexpr\linewidth-2\fboxsep\relax}%
   \setlength{\parskip}{4pt}\small
   {\bfseries\color{Accent}\addfontfeature{LetterSpace=6}\MakeUppercase{#1}}\par\smallskip}
  {\end{minipage}\end{lrbox}%
   \noindent\colorbox{Shade}{\usebox{\LPShadeBox}}\par\medskip}

% ---------- Answer areas ----------
% \begin{answer}[height] ... \end{answer}  boxed area; students write inside it
\newsavebox{\LPAnswerBox}
\newenvironment{answer}[1][3cm]
  {\par\medskip\begin{lrbox}{\LPAnswerBox}%
   \begin{minipage}[t][#1][t]{\dimexpr\linewidth-2\fboxsep-2\fboxrule\relax}%
   \setlength{\parskip}{4pt}%
   {\scriptsize\bfseries\color{Muted}\addfontfeature{LetterSpace=6}ANSWER}\par\smallskip}
  {\end{minipage}\end{lrbox}%
   \noindent\fcolorbox{Rule}{white}{\usebox{\LPAnswerBox}}\par\medskip}
% \answerlines[n]  n ruled writing lines (default 3)
\makeatletter
\newcommand{\answerlines}[1][3]{%
  \par\@tempcnta=#1\relax
  \@whilenum\@tempcnta>\z@\do{%
    \vspace{1.1em}\noindent{\color{Rule}\rule{\linewidth}{0.4pt}}\par
    \advance\@tempcnta\m@ne}}
\makeatother

% ---------- Tables ----------
\newcolumntype{C}{>{\centering\arraybackslash}X}       % centred, stretching column
\newcommand{\headcell}[1]{\textbf{\color{Accent}#1}}   % table column heading
\newcommand{\blankcell}{\rule[-0.8em]{0pt}{2.2em}}     % empty cell tall enough to write in
\newcommand{\tablecaption}[1]{\par\smallskip\noindent{\small\bfseries\color{Muted}#1}\par\nobreak\smallskip}

% ---------- Header block ----------
\newcommand{\infofield}[2]{%
  {\scriptsize\bfseries\color{Muted}\addfontfeature{LetterSpace=6}\MakeUppercase{#1}}\par\vspace{1pt}#2}
\newcommand{\worksheetheader}{%
  \thispagestyle{plain}%
  \noindent{\small\bfseries\color{Accent}\addfontfeature{LetterSpace=8}%
    \MakeUppercase{Assignment \assignmentnumber}\par}%
  \vspace{0.2em}%
  \noindent{\fontsize{23}{27}\selectfont\bfseries\color{Ink}\assignmenttitle\par}%
  \vspace{0.5em}%
  \noindent{\color{Ink}\rule{\linewidth}{0.8pt}}\par\vspace{0.2em}%
  \noindent\begin{tabularx}{\linewidth}{@{}>{\raggedright\arraybackslash}X>{\raggedright\arraybackslash}X@{}}
    \infofield{Student}{\studentname} & \infofield{Course}{\coursename} \\[0.7em]
    \infofield{Instructor}{\instructorname} & \infofield{Date}{\assignmentdate} \\
  \end{tabularx}%
  \par\vspace{0.1em}\noindent{\color{Rule}\rule{\linewidth}{0.4pt}}\par}

% ---------- Links (load last) ----------
\usepackage[colorlinks=true,linkcolor=Accent,citecolor=Accent,urlcolor=Accent]{hyperref}
\urlstyle{same}

% ---------- Assignment details ----------
% CUSTOMIZE: fill these in; they feed the header block and the running header.
\newcommand{\assignmentnumber}{1}
\newcommand{\assignmenttitle}{Calculus and Proof}
\newcommand{\studentname}{Student Name}
\newcommand{\coursename}{MATH 201 -- Calculus II}
\newcommand{\instructorname}{Dr.\ Instructor Name}
\newcommand{\assignmentdate}{\today}

\begin{document}

\worksheetheader

\begin{instructions}
  \begin{itemize}
    \item Show all work; answers without justification receive partial credit.
    \item Write final answers in the boxes provided and attach extra pages if needed.
    \item Calculators are not permitted.
  \end{itemize}
\end{instructions}

\section*{Part A -- Computation}

% ===== Problem with parts and nested sub-parts =====
\problem[10 pts]{Roots and Integrals}
Let $f(x) = x^{2} - 4x + 3$.
\begin{parts}
  \item Find all real roots of $f$.
  \item Evaluate the definite integral
    \[
      \int_{0}^{3} f(x)\,dx.
    \]
    \begin{parts}
      \item Find an antiderivative $F$ of $f$.
      \item Use $F$ to evaluate the integral.
    \end{parts}
  \item Explain in one sentence why your answer to (b) does not mean the region
    between the graph and the $x$-axis has zero area.
\end{parts}

\begin{solution}[Solution to (a)]
  Factoring gives $f(x) = (x-1)(x-3)$, so the roots are $x = 1$ and $x = 3$.
\end{solution}

\begin{answer}[3.5cm]
\end{answer}

% ===== Multiple choice =====
\problem[4 pts]{Multiple Choice}
\begin{note}
  Circle exactly one option for each question.
\end{note}
\begin{parts}
  \item What is the derivative of $\sin(x^{2})$?
    \begin{choices}
      \item $\cos(x^{2})$
      \item $2x\cos(x^{2})$
      \item $-2x\sin(x^{2})$
      \item $2x\sin(x^{2})$
    \end{choices}
  \item Which statement about the series $\displaystyle\sum_{n=1}^{\infty} \frac{1}{n^{2}}$ is true?
    \begin{choices}
      \item It diverges to infinity.
      \item It converges.
      \item It oscillates without a limit.
      \item Its convergence cannot be decided.
    \end{choices}
\end{parts}

\section*{Part B -- Reasoning}

% ===== Proof and explanation =====
\problem[6 pts]{Proof by Induction}
Prove that for every integer $n \ge 1$,
\begin{equation}\label{eq:odd-sum}
  1 + 3 + 5 + \cdots + (2n-1) = n^{2}.
\end{equation}
\begin{note}[Hint]
  Use induction on $n$.
\end{note}

\begin{proof}
  \emph{Base case.} For $n = 1$ the left side is $1 = 1^{2}$.

  \emph{Inductive step.} Assume \eqref{eq:odd-sum} holds for $n = k$. Then
  \begin{align*}
    1 + 3 + \cdots + (2k-1) + (2k+1) &= k^{2} + (2k+1) \\
                                     &= (k+1)^{2},
  \end{align*}
  so the statement holds for $n = k+1$. By induction it holds for all $n \ge 1$.
\end{proof}

\begin{solution}[Explanation]
  The identity says an $n \times n$ square can be built from L-shaped layers
  containing $1, 3, 5, \dots$ unit squares.
\end{solution}

\answerlines[3]

% ===== Table for answers or observations =====
\problem[5 pts]{Table of Values}
Complete the table for the function $f$ from Problem~1, then record your observations.

\tablecaption{Table 1. Values of $f(x) = x^{2} - 4x + 3$.}
\noindent
\begin{tabularx}{\linewidth}{@{}CCC@{}}
  \toprule
  \headcell{$x$} & \headcell{$f(x)$} & \headcell{Sign of $f(x)$} \\
  \midrule
  0 & 3 & $+$ \\
  1 & \blankcell & \blankcell \\
  2 & \blankcell & \blankcell \\
  3 & \blankcell & \blankcell \\
  4 & \blankcell & \blankcell \\
  \bottomrule
\end{tabularx}

\begin{answer}[3cm]
\end{answer}

% ---------------------------------------------------------------------------
% COPY-AND-EDIT SKELETONS: remove the leading % to use. Add as many as needed.
%
% \problem[5 pts]{Topic}
% Problem statement.
% \begin{parts}
%   \item First question.
%   \item Second question, with sub-questions:
%     \begin{parts}
%       \item First sub-question.
%       \item Second sub-question.
%     \end{parts}
% \end{parts}
% \begin{answer}[3cm]
% \end{answer}
%
% \problem[2 pts]{Multiple choice}
% Question text.
% \begin{choices}
%   \item First option
%   \item Second option
%   \item Third option
%   \item Fourth option
% \end{choices}
% ---------------------------------------------------------------------------

\end{document}
